\documentclass[10pt,a4paper]{amsart}
\usepackage[margin=19mm]{geometry}
\usepackage{amsmath,amssymb,mathtools}
\usepackage[scr=rsfs,cal=euler]{mathalfa}
\usepackage{xfrac}
\usepackage[unicode,hidelinks,bookmarks=false]{hyperref}
\newcommand{\ie}{i.e.\ }
\newcommand{\cf}{cf.\ }
\newcommand{\resp}{respectively\ }
\newcommand{\Subsection}{\subsection}
\providecommand{\nobreakdash}{\nobreak\hbox{-}}
\setlength{\emergencystretch}{2em}
\multlinegap0pt
\newtheorem{claim}{Claim}
\title{\textbf{TRAVELING WAVES IN reaction-diffusion EQUATIONS WITH General Delays}}
\author{William Barker, Nguyen Van Minh}
\date{Source: arXiv:2301.11504v3 (CC BY 4.0)}
\begin{document}
\maketitle

\noindent\emph{Source and licence.} \textbf{TRAVELING WAVES IN reaction-diffusion EQUATIONS WITH General Delays}. Version arXiv:2301.11504v3 (CC BY 4.0), distributed by arXiv under the Creative Commons Attribution 4.0 International licence, http://creativecommons.org/licenses/by/4.0/, as recorded in the arXiv licence metadata for this exact version and shown on the abstract page https://arxiv.org/abs/2301.11504v3. Full licence notice: \texttt{licenses/arxiv-2301.11504-CC-BY-4.0.txt}.\par\smallskip

\noindent\textit{Editorial scope.} Counted unit: the article's claim 6 (source0.tex lines 1030--1032). The article's complete original proof of it is delivered with it (source0.tex lines 1033--1171, one \texttt{proof} environment of 139 lines). No mathematics is written, repaired or re-derived: the statement and the proof are the article's own lines, in the article's own notation, and no retained line is substituted or re-flowed.  Retained uncounted as the source-local context the proof needs: lines 925--932; lines 956--958; lines 969--971.  Nothing was omitted from any retained span, so the package carries no omission record.  No retained line contains a citation, so the wrapper carries no bibliography and no bibliography entry.  No retained line contains a figure environment, an image-inclusion command or drawing code, so no image is delivered and the package declares no asset dependency.  Editorial formatting only, in addition: an AMS \texttt{amsart} preamble that restates the article's own declarations and macros used by the retained lines, namely the environments \texttt{claim} on separate counters, together with the standard packages those lines need.  The retained environments print in this document as Claim 1 in order of appearance, whereas the article prints the same items with the numbering of its own full document.  The complete native source of the article as distributed in the arXiv e-print for this exact version is bundled unchanged as \texttt{sources/136588/source0.tex}.
\begin{equation}\label{waveBZ}
\begin{cases}
\varphi_{1}^{\prime\prime}(t)-c\varphi_{1}^{\prime}(t+r_1)+\varphi_{1}(t+r_1)\left(
(1-r)-\varphi_{1}(t+r_1)-r\varphi_{2}(t+r_1-r_2)\right)  =0\\
\varphi_{2}^{\prime\prime}(t)-c\varphi_{2}^{\prime}(t+r_1)+b\varphi_{1}(t+r_1)\left(
1-\varphi_{2}(t+r_1)\right)  =0.
\end{cases},
\end{equation}

\begin{equation}\label{l1}
\lambda^2-c\lambda e^{r_1\lambda}+e^{r_1\lambda }=0.
\end{equation}

\begin{equation}\label{m1}
\mu^2-c\mu e^{r_1\mu}+be^{r_1\mu }=0.
\end{equation}

\begin{claim}\label{claim 6}
For sufficiently small $r_1$ and $c>2\sqrt{b}$, the vector function $(\varphi_1 ,\varphi_2)^T$ is a quasi-upper solution of (\ref{waveBZ}).
\end{claim}

\begin{proof}
{\bf Case $t\le - r_1 $}: 
Plugging these functions into the first equation in (\ref{waveBZ}) we have
for $t+r_1 \leq0$
\begin{align*}
&  \left[
\varphi_{1}^{\prime\prime}(t)-c\varphi_{1}^{\prime}(t+r_1)+\varphi
_{1}(t+r_1)\right]  -r\varphi_{1}(t+r_1)(1-\varphi_{2}(t+r_1-r_2))-\varphi_{1}^{2}(t+r_1)\\
&  =\left[  \frac{\lambda_{1}^{2}}{2}e^{\lambda_{1}t}-c\frac{\lambda_{1}} 
{2}e^{\lambda_{1}(t+r_1)}+\frac{1}{2}e^{\lambda_{1}(t+r_1)}\right]  -re^{\lambda_{1} 
(t+r_1)}\left(  1-\varphi_{2}(t+r_1-r_2)\right)  -\frac{1}{4}e^{2\lambda_{1}(t+r_1)}\\
&  =\left[  \lambda_{1}^{2}-c\lambda_{1}e^{r_1\lambda_1}+e^{r_1\lambda_1}\right]  \frac{1}{2}e^{\lambda_{1} 
t}-re^{\lambda_{1}(t+r_1)}\left(  1-\frac{1}{2} e^{\mu_1(t+r_1-r_2)}\right)  -\frac{1} 
{4}e^{2\lambda_{1}(t+r_1)}\\
&  =-re^{\lambda_{1}(t+r_1)}\left(  1-\frac{1}{2} e^{\mu_1(t+r_1-r_2)}\right)  -\frac{1} 
{4}e^{2\lambda_{1}(t+r_1)}
\leq 0.
\end{align*}

\medskip
For the second equation in (\ref{waveBZ}), for $t\le -r_1$ we have 
\begin{align*}
&  \varphi_{2}^{\prime\prime}(t)-c\varphi_{2}^{\prime}(t+r_1)+b\varphi
_{1}(t+r_1)\left(  1-\varphi_{2}(t+r_1)\right)  \\
&
=\frac{\mu_{1}^{2}}{2}e^{\mu_{1}t}-c\frac{\mu_{1}}{2}e^{\mu_{1}(t+r_1)}+\frac
{b}{2}e^{\lambda_{1}(t+r_1)}\left(  1-\frac{1}{2}e^{\mu_{1}(t+r_1)}\right)  \\
&  =\frac{1}{2}\left( \mu^2_1-c\mu_1e^{r_1\mu_1}+be^{r_1\mu_1} \right)  e^{\mu_1t} 
+\frac{b}{2} \left(  e^{\lambda_1(t+r_1)} -e^{\mu_1(t+r_1)}     \right)
-\frac{b}{4} e^{(\lambda_1+\mu_1)(t+r_1)}\\
& = \frac{b}{2} \left(  e^{\lambda_1(t+r_1)} -e^{\mu_1(t+r_1)}     \right)
-\frac{b}{4} e^{(\lambda_1+\mu_1)(t+r_1)} \le 0
\end{align*}
since $\lambda_1 >\mu_1$ and $t+r_1 \le 0$.

\medskip\noindent
{\bf Case: $-r_1 \le t\le 0$}:
\begin{align*}
A:=&  \left[
\varphi_{1}^{\prime\prime}(t)-c\varphi_{1}^{\prime}(t+r_1)+\varphi
_{1}(t+r_1)\right]  -r\varphi_{1}(t+r_1)(1-\varphi_{2}(t+r_1-r_2))-\varphi_{1}^{2}(t+r_1)\\
&  =\left[  \frac{\lambda_{1}^{2}}{2}e^{\lambda_{1}t}-c\frac{\lambda_{1}} 
{2}e^{-\lambda_{1}(t+r_1)}+1-\frac{1}{2}e^{-\lambda_{1}(t+r_1)}\right]  \\
& \hspace{.5cm} -r\left( 1-\frac{1}{2}e^{-\lambda_{1}(t+r_1)}\right) \left(  1-\varphi_{2}(t+r_1-r_2)\right)  -\varphi_{1}^{2}(t+r_1)\\
\end{align*}
As $\lambda_1$ is a root of Eq.(\ref{l1})
\begin{align*}
A&= \frac{c\lambda_1}{2}\left(e^{\lambda_1(t+r_1)}-e^{-\lambda_1(t+r_1)} \right) + 1-\frac{1}{2}\left( e^{\lambda_1(t+r_1)}+e^{-\lambda_1(t+r_1)} \right) \\
& \hspace{.5cm} -r\left( 1-\frac{1}{2}e^{-\lambda_{1}(t+r_1)}\right) \left(  1-\varphi_{2}(t+r_1-r_2)\right)  -\varphi_{1}^{2}(t+r_1)\\
&  =c\lambda_1 \sinh (\lambda_1(t+r_1))+1-\cosh (\lambda_1(t+r_1)) \\
& \hspace{.5cm} -r\left( 1-\frac{1}{2}e^{-\lambda_{1}(t+r_1)}\right) \left(  1-\varphi_{2}(t+r_1-r_2)\right)  -\varphi_{1}^{2}(t+r_1).
\end{align*}
As $-r_1\le t\le 0$ and $r_1$ is sufficiently small, using the power expansions of $\sinh x$ and $\cosh x$ we have
$$
\sinh (\lambda_1(t+r_1)) \approx \lambda_1(t+r_1 )+o(r^2_1), \quad 1-\cosh (\lambda_1(t+r_1))\approx - \frac{(\lambda_1(t+r_1))^2}{2}+o(r^2_1).
$$
Finally, we have 
\begin{align*}
A&  = c\lambda_1^2(t+r_1) - \frac{(\lambda_1(t+r_1))^2}{2}+o(r^2_1) \\
&\hspace{.5cm }  -r\left( 1-\frac{1}{2}e^{-\lambda_{1}(t+r_1)}\right) \left(  1-\varphi_{2}(t+r_1-r_2)\right)  -\varphi_{1}^{2}(t+r_1)\le 0,
\end{align*}
because when $r_1$ is sufficiently small, the limit of the right hand side is $-(r+1)/4$ as $r_1\to 0$.

\medskip
For the second equation we have
\begin{align*}
B&:= \varphi_{2}^{\prime\prime}(t)-c\varphi_{2}^{\prime}(t+r_1)+b\varphi_{1}(t+r_1)\left(
1-\varphi_{2}(t+r_1)\right)  \\
&= \frac{\mu_1^2}{2}e^{\mu_1t}-\frac{c\mu_1}{2}e^{-\mu_1(t+r_1)}+b\left( 1-\frac{1}{2}e^{-\lambda_1(t+r_1)}\right) \frac{1}{2}e^{-\mu_1(t+r_1)}\\
&= \left[ \frac{\mu_1^2}{2}e^{\mu_1t} -\frac{c\mu_1}{2}e^{\mu_1(t+r_1)} +\frac{b}{2}e^{\mu_1(t+r_1)} \right] 
+\frac{c\mu_1}{2}e^{\mu_1(t+r_1)}  -\frac{c\mu_1}{2}e^{-\mu_1(t+r_1)} \\
&\hspace{0.5cm} - \frac{b}{2}e^{\mu_1(t+r_1)}  +\frac{b}{2}e^{-\mu_1(t+r_1)} -\frac{b}{4}e^{-(\lambda_1+\mu_1)(t+r_1)} .
\end{align*}
As $\mu_1$ is a root of Eq.(\ref{m1}), 
\begin{align*}
B&=
\frac{c\mu_1}{2}e^{\mu_1(t+r_1)}  -\frac{c\mu_1}{2}e^{-\mu_1(t+r_1)}  - \frac{b}{2}e^{\mu_1(t+r_1)}  +\frac{b}{2}e^{-\mu_1(t+r_1)} -\frac{b}{4}e^{-(\lambda_1+\mu_1)(t+r_1)} .
\end{align*}
As $-r_1\le t\le 0$ and the limit of the right hand side is $-b/4<0$ as $r_1\to 0$, it follows that for sufficiently small $r_1$ we will have  $B\le 0$.


\medskip
{\bf Case  $t>0$}:
\begin{align*}
A&:= \left[
\varphi_{1}^{\prime\prime}(t)-c\varphi_{1}^{\prime}(t+r_1)+\varphi
_{1}(t+r_1)\right]  -r\varphi_{1}(t+r_1)(1-\varphi_{2}(t+r_1-r_2))-\varphi_{1}^{2}(t+r_1)\\
&  = \left[  \frac{-\lambda_{1}^{2}}{2}e^{-\lambda_{1}t}-c\frac{\lambda_{1}} 
{2}e^{-\lambda_{1}(t+r_1)}+1-\frac{1}{2}e^{-\lambda_{1}(t+r_1)}\right]  \\
& \hspace{.5cm}-r\varphi_{1}(t+r_1)(1-\varphi_{2}(t+r_1-r_2))-\varphi_{1}^{2}(t+r_1)\\
 &=  \left[  \frac{-\lambda_{1}^{2}}{2}e^{-\lambda_{1}t}+c\frac{\lambda_{1}} 
{2}e^{-\lambda_{1}t+\lambda_1 r_1}-\frac{1}{2}e^{-\lambda_{1}t+\lambda_1r_1}\right] \\
& -c\frac{\lambda_{1}} 
{2}e^{-\lambda_{1}t+\lambda_1 r_1} -c\frac{\lambda_{1}} 
{2}e^{-\lambda_{1}t-\lambda_1 r_1} +\frac{1}{2}e^{-\lambda_{1}t+\lambda_1 r_1} +1-\frac{1}{2}e^{-\lambda_{1}(t+r_1)} \\
&-r\varphi_{1}(t+r_1)(1-\varphi_{2}(t+r_1-r_2))-\varphi_{1}^{2}(t+r_1) 
 \end{align*}
Since $\lambda_1$ is a root of Eq.(\ref{l1}), we have
\begin{align*}
&  \frac{\lambda_{1}^{2}}{2}e^{-\lambda_{1}t}-c\frac{\lambda_{1}} 
{2}e^{-\lambda_{1}t+\lambda_1 r_1}+\frac{1}{2}e^{-\lambda_{1}t+\lambda_1r_1}=0.
\end{align*}
Hence,
\begin{align*}
A&= -c\frac{\lambda_{1}} 
{2}e^{-\lambda_{1}t+\lambda_1 r_1} -c\frac{\lambda_{1}} 
{2}e^{-\lambda_{1}t-\lambda_1 r_1} +\frac{1}{2}e^{-\lambda_{1}t+\lambda_1 r_1} +1-\frac{1}{2}e^{-\lambda_{1}(t+r_1)} \\
&-r\varphi_{1}(t+r_1)(1-\varphi_{2}(t+r_1-r_2))-\varphi_{1}^{2}(t+r_1) \\
&= \left[ 1-c\lambda_1 \cosh(\lambda_1r_1)+\sinh(\lambda_1r_1)\right] e^{-\lambda_1 t} \\
& -r\varphi_{1}(t+r_1)(1-\varphi_{2}(t+r_1-r_2))-\varphi_{1}^{2}(t+r_1).
\end{align*}
Since
$$
\lim_{r_1\downarrow 0} \left[ 1-c\lambda_1 \cosh(\lambda_1r_1)+\sinh(\lambda_1r_1)\right] <0,
$$
it follows that $A\le 0$ for sufficiently small $r_1$.

For the second equation
\begin{align*}
B&:= \varphi_{2}^{\prime\prime}(t)-c\varphi_{2}^{\prime}(t+r_1)+b\varphi_{1}(t+r_1)\left(
1-\varphi_{2}(t+r_1)\right)  \\
&= -\frac{\mu_1^2}{2}e^{-\mu_1t}-\frac{c\mu_1}{2}e^{-\mu_1(t+r_1)}+b\left( 1-\frac{1}{2}e^{-\lambda_1(t+r_1)}\right) \frac{1}{2}e^{-\mu_1(t+r_1)}\\
&= \left[  -\frac{\mu_1^2}{2}e^{-\mu_1t} +\frac{\mu_1}{2}e^{-\mu_1t+\mu r_1} -\frac{b}{2}e^{-\mu_1 t+\mu_1r_1)} \right] 
\\
& \hspace{0.5cm} -\frac{c\mu_1}{2}e^{-\mu_1t+\mu r_1} -\frac{c\mu_1}{2}e^{-\mu_1t-\mu r_1} +\frac{b}{2}e^{-\mu_1 t+\mu_1r_1)} \\
& \hspace{0.5cm}+ \frac{b}{2}e^{-\mu_1(t+r_1)} -\frac{b}{4}e^{-(\lambda_1+\mu_1)(t+r_1)} .
\end{align*}
As $\mu_1$ is a root of Eq.(\ref{m1}),
\begin{align*}
B&= -\frac{c\mu_1}{2}e^{-\mu_1t+\mu_1 r_1} -\frac{c\mu_1}{2}e^{-\mu_1t-\mu_1 r_1} +\frac{b}{2}e^{-\mu_1 t+\mu_1r_1} \\
& \hspace{0.5cm} +\frac{b}{2}e^{-\mu_1(t+r_1)} -\frac{b}{4}e^{-(\lambda_1+\mu_1)(t+r_1)} \\
& = \left(\frac{b}{2}-c\mu_1 \right)\cosh (\mu_1 r_1)e^{-\mu_1t} -\frac{b}{4}e^{-(\lambda_1+\mu_1)(t+r_1)} .
\end{align*}
Since $c>2\sqrt{b}$ we have $c^2>4b$, and 
$$
\frac{b}{2} < \frac{c^2}{8} < \frac{c(c+\sqrt{c^2-4b})}{2}=c\mu_1 .
$$
This yields  $B\le 0.$ The claim is proved.
\end{proof}

\end{document}
